\documentclass[11pt]{article}
\usepackage[T1]{fontenc}
\usepackage[utf8]{inputenc}
\usepackage{lmodern}
\usepackage[margin=1in]{geometry}
\usepackage{amsmath,amssymb,amsthm,mathtools,microtype,booktabs}
\usepackage[hidelinks]{hyperref}
\hypersetup{pdftitle={The Reset Calculus: Antimatroids and Forest Models},
  pdfauthor={KFOL research working note}}
\newtheorem{theorem}{Theorem}[section]
\newtheorem{lemma}[theorem]{Lemma}
\newtheorem{corollary}[theorem]{Corollary}
\newtheorem{proposition}[theorem]{Proposition}
\newcommand{\Pow}{\mathcal P_+}
\newcommand{\del}{\mathbin{\triangleleft}}
\newcommand{\bas}{\operatorname{bas}}
\title{The Reset Calculus:\\Antimatroids and Forest Models}
\author{KFOL research working note}
\date{30 October 2026}
\begin{document}
\maketitle
\begin{abstract}
We isolate one operation derived from rational owned choice: retain a
resolved right operand, otherwise take union. It generates exactly the
binary-generated part of the idempotent, atom-preserving fragment;
ternary operations already go beyond it. Its complete term clone,
including mixed resolved and ambiguous inputs, is represented by
antimatroids with a unique initial element. Canonical prerequisite
antichains give polynomial algorithms for equality, realizability,
counterexamples, and synthesis. Substitution has an explicit formula.
All numbers of players at least three give the same equational theory,
generated by a three-element chain. This generator is exactly the dual
of Waldhauser's known groupoid $G_3$, and its clone is minimal: every
nonprojection term recovers reset by identifying variables. The first
five free-algebra sizes are $1,4,18,140,2980$. Among conservative algebras,
two equations characterize its entire theory: the models are exactly
forest orders. Comb-shaped forests give subdirectly irreducible models
of every infinite cardinality, so the variety is residually large.
Mixed-input semantics, three-element completeness, reconstruction of
forest orders, and the generic comb-monolith argument are kernel-checked
in Lean. The antimatroid correspondence and membership of forest models
in the reset variety have written proofs and independent finite checks.
The small generator and dependency structures with choice are established
prior art, used to describe the fragment and its models.
\end{abstract}

\section{Why study reset?}

Fix a winner set $I$, $|I|=K\geq3$. Values are nonempty subsets of $I$.
The public rational owned-choice operation is
\[
 A\star_r B=
 \begin{cases}\{r\},&A=\{r\}\text{ or }B=\{r\},\\
 A\cup B,&\text{otherwise}.
 \end{cases}
\]
It composes the possible winners of full pure subgame-perfect
equilibria in finite public games with unique-winner leaves. The
recurrence and the full-profile game semantics are proved in
\texttt{Basic.lean}. Three owners define ordinary union. Starting with
$A\cup B$, suppress once by each owner, replacing the current $U$ by
$U\star_r B$. The resulting operation is
\begin{equation}\label{eq:reset}
 A\del B=
 \begin{cases}
 B,&|B|=1,\\
 A\cup B,&|B|\geq2.
 \end{cases}
\end{equation}
The order of suppression does not matter. The relative-reset laws and
their expansion into owned-choice syntax are kernel-checked.

Study the single-operation language
\[
 t::=x_j\mid t\del t.
\]
We call this the \emph{reset calculus}. It has no winner constants or
internal role switches. Its operation commutes with every permutation
of winner roles. The question is whether a fragment this small still
has a useful exact account of composition.

The earlier \texttt{general.tex} identified a right-normal-band quotient
by observing only all-singleton and all-ambiguous valuations. That
quotient loses mixed behavior:
\[
 ((\{c\}\del\{a\})\del\{a,b\})=\{a,b\},\qquad
 \{c\}\del(\{a\}\del\{a,b\})=\{a,b,c\}.
\]
The representation below keeps that distinction.

Reset is a proper fragment of owned choice. Every reset term returns a
singleton on singleton inputs. An owner choice between two distinct
other winners returns their pair. Neither that choice nor ordinary
union can therefore be expressed by reset alone.

\section{An intrinsic binary fragment}

The operation also has a characterization inside owned choice.
A term \emph{preserves atoms} if it returns a singleton whenever all
inputs are singletons. Work first with constant-free, switch-free
owned-choice terms, which are automatically idempotent.

\begin{lemma}[One-constant rigidity]\label{lem:rigidity}
A unary owned-choice expression whose only constant is $\{r\}$ is
either constantly $\{r\}$ or inflationary: it contains its input on
every valuation. Consequently, returning $\{r\}$ at any other singleton
$\{q\}$ already forces it to be constant.
\end{lemma}
\begin{proof}
Induct on expressions. Combining two inflationary expressions preserves
inclusion of the input: if an owner wins outright, that singleton
already contains the input. Combining an inflationary expression with
the constant $\{r\}$ returns the constant when the owner is $r$;
otherwise it remains inflationary, because the constant operand cannot
give a different owner an outright win. Two constant expressions remain
constant. An inflationary function cannot send $\{q\}$ to $\{r\}$ when
$q\ne r$. The dichotomy and this rigidity consequence are
kernel-checked in \texttt{OneConstant.lean}.
\end{proof}

\begin{theorem}[Exact binary atom-preserving fragment]\label{thm:binary}
For $K\geq3$, the binary atom-preserving operations definable by
constant-free owned choices are exactly
\[
 x,\qquad y,\qquad x\del y,\qquad y\del x.
\]
Hence the reset clone is exactly the part generated by the binary
operations of the idempotent, atom-preserving owned-choice clone.
\end{theorem}
\begin{proof}
The atomic restriction of an idempotent atom-preserving owned term is
a projection. This is the established $S_2$ application in
\texttt{general.tex}: $\operatorname{Pol}(\iota_3)=S_2$
\cite[Proposition~2.2]{Szendrei2012}, and an idempotent operation on
$K\geq3$ atoms has image size at least three, so the essentially unary
case of $S_2$ must apply.

Suppose the atomic restriction projects $x$. If $y$ does not occur,
idempotency makes the whole term $x$. Otherwise both variables occur.
Freeze $x=\{r\}$, leaving a unary expression in $y$ with only that
constant. Its atomic outputs are $\{r\}$, so
Lemma~\ref{lem:rigidity} makes it constant on all $y$.

Now let $x$ be ambiguous. If $y$ is ambiguous too, every node computes
union, giving $x\cup y$. If $y=\{r\}$, specialize instead in $x$.
Its atomic restriction is the identity, so rigidity makes it
inflationary. Its result is contained in $x\cup\{r\}$, since owned
choice never introduces a winner absent from its inputs. It contains
$r$ as well: under this valuation every singleton subexpression is
$\{r\}$, by induction, so no different owner can suppress the
contribution of an occurring $y$. The value is therefore
$x\cup\{r\}$. These cases give exactly $y\del x$.
The opposite projection gives the other two operations, all of which
are definable.
\end{proof}

The earlier idempotent-elimination theorem in \texttt{general.tex}
extends this characterization to idempotent full-signature terms,
including constants and internal switches in their original syntax.
The binary theorem uses that earlier mathematical theorem and the
external atomic-clone classification; its complete proof is not yet a
Lean theorem.

The binary restriction is strict. For an owner $a$, the ternary operation
\[
 H_a(x,y,z)=(y\star_a z)\del x
\]
is idempotent and atom-preserving. Set $x=\{a,b\}$ and $z=\{c\}$.
At $y=\{a\}$ its value omits $c$, while at $y=\{b\}$ it contains $c$.
The ambiguity statuses are unchanged, and only $z$ supplies $c$.
The selector semantics below therefore rules out a reset expression
for $H_a$. Thus reset is an intrinsic binary core, while higher-arity
atom-preserving operations still retain owner-sensitive behavior.

\section{Exact prerequisite semantics}

Let $N$ be the finite set of input positions, including any declared
unused positions. For a valuation $S=(S_j)_{j\in N}$ put
$B=\{j:|S_j|\geq2\}$. This is a set of input positions, not winner roles.
The root $\rho(t)$ is its rightmost variable:
\[
 \rho(x_j)=j,\qquad \rho(s\del t)=\rho(t).
\]

\begin{lemma}[Resolution and selection]\label{lem:selection}
A reset term is ambiguous exactly when its root input is ambiguous.
There is a selector $E_t:2^N\to2^N$ with
\begin{align*}
 E_{x_j}(B)&=\{j\},\\
 E_{s\del t}(B)&=E_t(B)\cup
 \begin{cases}E_s(B),&\rho(t)\in B,\\\varnothing,&\rho(t)\notin B,\end{cases}\\
 t(S)&=\bigcup_{j\in E_t(B)}S_j.
\end{align*}
Every selector is monotone, always selects its root, and selecting
position $j$ is independent of whether $j$ itself belongs to $B$.
\end{lemma}
\begin{proof}
If the right child is resolved, reset returns it. If it is ambiguous,
reset unions the children and remains ambiguous. This proves the first
three statements by induction. Monotonicity follows from the positive
recursion. For self-independence, at a node with root $j$ that input is
always selected. Otherwise its new guard concerns a different position,
and both child selection predicates are independent of $j$ by induction.
\end{proof}

For each position $j$, record the inclusion-minimal sets $A\subseteq N$
such that $j\in E_t(A)$. Denote this antichain by $\mathcal G_j(t)$.
An empty antichain means the position is unused. Then
\begin{equation}\label{eq:guards}
 j\in E_t(B)\quad\Longleftrightarrow\quad
       \exists A\in\mathcal G_j(t)\; A\subseteq B.
\end{equation}
For the root its antichain is $\{\varnothing\}$, and every other
antichain omits $\varnothing$. No minimal prerequisite for $j$ contains
$j$, by self-independence.

These tables are computed directly. A variable contributes its empty
guard. At $s\del t$, retain the right table and add $\rho(t)$ to each
left guard, then delete duplicates and every guard containing a smaller
one for the same output. Before minimization there is at most one guard
per variable occurrence. A tree traversal can collect these guards and
normalize them in polynomial time.

\begin{theorem}[Canonical equality and three-element completeness]
\label{thm:three}
Two reset terms are equal on all nonempty-set valuations if and only if
their prerequisite tables agree. The result is independent of $I$ for
every $|I|\geq3$. It suffices to evaluate on a fixed three-element chain
\[
 Z=\{a\}\subset U=\{a,b\}\subset T=\{a,b,c\}.
\]
Consequently equality is decidable in polynomial time in reset-tree size.
A failed equality has a polynomially constructible countervaluation
using only $Z,U,T$.
\end{theorem}
\begin{proof}
Equal tables give equal selectors and therefore equal values.
If two tables differ, some selector inclusion differs at $B,j$.
Assign $T$ to position $j$, $U$ to the other positions in $B$, and $Z$
elsewhere. Setting $j$ to ambiguous does not change whether $j$ is
selected, by Lemma~\ref{lem:selection}. The marker $c$ occurs in the
result exactly when that position is selected. Thus the terms differ.
Distinct minimal antichains give such $B,j$: take a minimal prerequisite
on one side not containing a prerequisite from the other. All operations
on finite guard tables are polynomial.
\end{proof}

The chain is closed under reset. Numbering it $0<1<2$, its table is
\begin{equation}\label{eq:three-table}
\begin{array}{c|ccc}
\del&0&1&2\\\hline
0&0&1&2\\
1&0&1&2\\
2&0&2&2
\end{array}
\end{equation}
This is the dual of $G_3$ in Waldhauser's classification
\cite[Theorem~3.4, table $G_3$]{Waldhauser2006}, under
$0\mapsto b$, $1\mapsto a$, $2\mapsto c$.
This identifies the reset variety with the variety generated by that
established three-element algebra.

\begin{corollary}[Minimum finite generator]
The three-element algebra generates the common variety of all reset
algebras on $\Pow(I)$ for $|I|\geq3$, and no smaller algebra generates
that variety.
\end{corollary}
\begin{proof}
Theorem~\ref{thm:three} identifies their identities. The three-element
algebra is nonassociative:
$(2\del0)\del1=1$, while $2\del(0\del1)=2$.
Every idempotent binary operation on two elements is one of the two
projections or two semilattice operations; all four are associative.
Any generator would have to satisfy the valid identity $x\del x=x$,
so a smaller generator is impossible.
\end{proof}

Even a two-value input domain that is not closed under reset cannot
test all identities: associativity holds whenever its three arguments
use at most two distinct winner sets. A failure requires the middle
argument to be atomic and the last ambiguous, with a contribution of
the first lying outside their union. Two distinct input values cannot
supply that contribution.

\section{The antimatroid represented by a term}

Prerequisites can describe alternative routes. For example, let
\[
 t=((x_3\del x_1)\del x_0)\del((x_3\del x_2)\del x_0).
\]
Its canonical table is
\[
 \mathcal G_0=\{\varnothing\},\quad
 \mathcal G_1=\mathcal G_2=\{\{0\}\},\quad
 \mathcal G_3=\{\{0,1\},\{0,2\}\}.
\]
Input $0$ is always retained. Inputs $1,2$ are retained when $0$ is
ambiguous, while input $3$ additionally requires either $1$ or $2$ to
be ambiguous. This alternative cannot be described by a single
partial order of prerequisites. Feasible sets record which positions
can be reached by following the alternatives.

We use the standard finite feasible-set definition
\cite[Section~7]{Czedli2014}. An antimatroid on its actual support $X$
is a union-closed family $\mathcal F\subseteq2^X$ containing
$\varnothing$ and $X$, with the accessibility property: every nonempty
feasible set can lose one element and remain feasible. We require
exactly one feasible singleton $\{r\}$. Accessibility then makes $r$
belong to every nonempty feasible set. Unused input positions lie
outside $X$.

Write
\[
 \bas_{\mathcal F}(B)=\bigcup\{A\in\mathcal F:A\subseteq B\},\qquad
 e_{\mathcal F}(B)=\{j:j\in\bas_{\mathcal F}(B\cup\{j\})\}.
\]
The basis is the largest feasible subset of $B$. The set $e_{\mathcal F}(B)$
consists of that basis together with its available one-element
extensions. To verify this description, take a feasible subset of
$B\cup\{j\}$ containing $j$ and an accessible order of it. Its prefix
before $j$ is feasible inside $B$; unioning with the basis makes $j$
an available extension. The converse is immediate.

\begin{theorem}[Exact antimatroid representation]\label{thm:anti}
Reset term functions on a fixed input set $N$ correspond bijectively to
antimatroids with one feasible singleton on nonempty supports
$X\subseteq N$. The correspondence is
\[
 E_t=e_{\mathcal F_t},\qquad
 \mathcal F_t=\{A\subseteq X:A\subseteq E_t(A)\}.
\]
For reset itself the family operation is
\begin{equation}\label{eq:family-product}
 \mathcal F\del\mathcal G=
 \{\varnothing\}\cup
 \{A\cup C:A\in\mathcal F,\ C\in\mathcal G\setminus\{\varnothing\}\}.
\end{equation}
The root of the product is the root of $\mathcal G$.
\end{theorem}
\begin{proof}
The variable $x_j$ gives $\{\varnothing,\{j\}\}$.
The family in \eqref{eq:family-product} is union-closed. To order a
nonempty member $A\cup C$ accessibly, first take an accessible order of
$C$ in $\mathcal G$, then one of $A$ in $\mathcal F$, skipping positions
already included. Every prefix is a union of a feasible set of
$\mathcal F$ with a nonempty feasible set of $\mathcal G$. This proves
accessibility, full support, and the unique initial element.

If $r$ is the root of $\mathcal G$, its basis obeys
\[
 \bas_{\mathcal F\del\mathcal G}(B)=
 \begin{cases}
 \varnothing,&r\notin B,\\
 \bas_{\mathcal F}(B)\cup\bas_{\mathcal G}(B),&r\in B.
 \end{cases}
\]
Applying the definition of $e$ gives exactly the selector recursion
in Lemma~\ref{lem:selection}. Also
$A\subseteq e_{\mathcal F}(A)$ iff $A$ is feasible: membership of
$j\in A$ in this readiness set means $j\in\bas_{\mathcal F}(A)$.
This proves faithful representation for every term.

Conversely, an accessible order $r,j_2,\ldots,j_m$ of the full support
is a basic word. Its path term
\[
 x_{j_m}\del x_{j_{m-1}}\del\cdots\del x_r
\]
with \emph{left} parentheses has the feasible family of its prefixes.
Here the root $r$ is the last argument of the term, while it is the
first element of the accessible order. Combining two terms with the
same root by reset is pointwise union. On families it is union closure
of the two families: each already contains the common singleton root.
Every feasible set extends to a basic word, by accessibility and union
closure. Thus joining the path terms for all basic words gives exactly
the prescribed antimatroid. This is an application of the standard
basic-word representation, with reset supplying its composition.
\end{proof}

The parentheses in the path display are essential: for an order
$r,a,b$, the term is $(x_b\del x_a)\del x_r$. It selects the initial
segment through the first resolved input. Repeated labels in more
general terms are handled by \eqref{eq:family-product}; replacing the
occurrence tree by a graph on variable names can incorrectly combine
paths from different occurrences.

\section{Polynomial recognition and synthesis}

The family representation can be exponentially large. Prerequisite
antichains retain a polynomial-size description for a given term.
They also have an exact realizability criterion.

Given a candidate root $r$ and minimal antichains $\mathcal G_j$,
call a set $A$ \emph{constructible} if its positions can be ordered
$a_1,\ldots,a_m$ such that for every $\ell$ some member of
$\mathcal G_{a_\ell}$ lies in $\{a_1,\ldots,a_{\ell-1}\}$.
This is checked greedily: start empty and add any enabled member of
$A$. Enabling is monotone, so a greedy choice cannot prevent success.

\begin{theorem}[Exact realizability]\label{thm:synthesis}
A table is the canonical table of a reset term exactly when:
\begin{enumerate}
\item $\mathcal G_r=\{\varnothing\}$ and no other antichain contains
$\varnothing$;
\item every member of $\mathcal G_j$ omits $j$;
\item every member of every antichain is constructible.
\end{enumerate}
Recognition is polynomial in the table description. A realizable
table with $M$ prerequisites on $n$ input positions has a term
realization of size $O(nM)$.
\end{theorem}
\begin{proof}
For necessity, roots and self-independence give the first two
conditions. A minimal prerequisite $A$ enabling $j$ is feasible in the
represented antimatroid: readiness at $A$ has a feasible predecessor
inside $A$, and minimality makes that predecessor all of $A$.
Any accessible order of $A$ verifies constructibility.

For sufficiency, for each prerequisite $A\in\mathcal G_j$ choose its
greedy order $a_1,\ldots,a_m$ and make the path term with order
$a_1,\ldots,a_m,j$. Its last position is not repeated, by condition 2.
The first is the unique unconditionally enabled position $r$.
The endpoint has exactly prerequisite $A$. Every earlier position on
the path has a prerequisite contained in the preceding prefix, by
constructibility. Thus all the path's extra contributions are already
permitted by the candidate table. Join these paths by reset with their
common root. Every required prerequisite is present and no unwanted
contribution is introduced. There are $M$ paths of length at most
$n+1$.
\end{proof}

The constructibility condition has a simpler equivalent local form:
\begin{equation}\label{eq:complete}
 A\in\mathcal G_j,\ k\in A
 \quad\Longrightarrow\quad
 \exists B\in\mathcal G_k\; B\subseteq A.
\end{equation}
Constructibility implies this completeness condition. Conversely,
induct on the size of a prerequisite $A$. For each $k\in A$, choose
$B\in\mathcal G_k$ inside $A$. Since $k\notin B$, it is smaller than
$A$ and constructible by induction. Append $k$ to its construction.
Unions of constructible sets are constructible: concatenate their
orders, skipping repeated positions. Unioning these constructions
reaches all of $A$.

On actual support, the antichain, nonemptiness, self-exclusion, and
completeness conditions are precisely the axioms for a \emph{dependency
structure with choice} (DSC) in
\cite[Definition~2.9]{BazermanMinichielloPuzio2023}. Their correspondence
with antimatroids is established in that work
\cite[Theorem~3.15]{BazermanMinichielloPuzio2023}. Thus the prerequisite
model reuses DSCs, with one unconditionally enabled event, as a concrete
presentation of reset term functions. The implementation checks the
local condition and uses greedy orders for synthesis.

An exponential antimatroid is therefore not needed for either
equality or synthesis. These are algorithms for formulas over the
primitive reset operation. Its expansion into owned choices may
duplicate subterms; no polynomial bound on that separate expansion
is claimed.

\section{Substitution, counting, and minimality}

Let an outer table have root $r$ and prerequisites $\mathcal G_i$.
Substitute inner tables with roots $r_i$ and prerequisites
$\mathcal H^i_j$. Since inner ambiguity is determined exactly by its
root, substitution gives
\begin{equation}\label{eq:substitution}
 r_{\rm new}=r_r,\qquad
 \mathcal G^{\rm new}_j=
 \min\bigcup_i
 \left\{\{r_k:k\in A\}\cup C:
       A\in\mathcal G_i,\ C\in\mathcal H^i_j\right\}.
\end{equation}
Here $\min$ deletes nonminimal sets. The formula follows by substituting
the inner output unions into the outer output union, and substituting
each inner root status into an outer guard. It is polynomial in the
two explicit table descriptions. It respects actual term substitution,
including mixed inputs, rather than only a homogeneous observation.

Removing the unique initial element from every nonempty feasible set
leaves an arbitrary antimatroid on the remaining support. Conversely,
adjoining one initial element to every feasible set gives the unique-root
form. If $a_m$ counts labeled antimatroids on a full $m$-element support,
with $a_0=1$, the number of $n$-ary reset term functions is therefore
\[
 f_n=\sum_{m=1}^n\binom nm\,m\,a_{m-1}.
\]
Independent finite enumerations give
\[
 (a_0,\ldots,a_4)=(1,1,3,22,485),\qquad
 (f_1,\ldots,f_5)=(1,4,18,140,2980).
\]
The antimatroid counts are counts of known combinatorial objects;
the equation describes their exact relationship to reset operations.

\begin{theorem}[Identification-minor minimality]\label{thm:minimal}
Every nonprojection reset term has a binary identification minor equal
to reset. Hence its term clone is minimal.
\end{theorem}
\begin{proof}
A term selects all its occurring variables when every input is
ambiguous. Each of them is essential, as witnessed by assigning a proper
ambiguous set to all inputs except one assigned a larger set.
A nonprojection therefore has support of size at least two. Identify
its root variable with $y$ and every other occurring variable with $x$.
The resulting binary term has root $y$ and both arguments essential.
Its $x$ prerequisite cannot contain $x$, cannot be empty, and can only
use $y$. Thus it is $\{\{y\}\}$, while the $y$ prerequisite is empty.
This is exactly $x\del y$. Minimality follows.
\end{proof}

Minimality is already known for $G_3$ and its dual in
Waldhauser's work. The direct proof here records the stronger explicit
recovery by input identification, using the prerequisite representation.

\section{Two equations and the conservative models}

An algebra $(A,\del)$ is \emph{conservative} if $x\del y\in\{x,y\}$.
The three-element generator is conservative, although winner-set reset
and free reset algebras generally are not. Here we classify conservative
models of the equational theory; their elements need not themselves be
concrete winner sets.

A \emph{forest order} is a partial order in which the upper set of each
element is linearly ordered. In the finite case its Hasse diagram is a
rooted forest, with roots above their descendants. Define
\begin{equation}\label{eq:forest-op}
 x\del_P y=\begin{cases}x,&y\leq_P x,\\y,&\text{otherwise}.\end{cases}
\end{equation}
Comparable elements yield the higher one; incomparable elements yield
the right argument.

\begin{theorem}[Conservative presentation]\label{thm:forest}
For a conservative algebra, the following are equivalent:
\begin{enumerate}
\item it satisfies every reset identity;
\item it satisfies the two equations
\begin{align}
 (x\del y)\del x&=y\del x,\label{eq:absorb}\\
 x\del(y\del z)&=y\del(x\del z);\label{eq:left-comm}
\end{align}
\item it has form \eqref{eq:forest-op} for a unique forest order.
\end{enumerate}
The order is recovered by
$x\leq_P y$ iff $x\del y=y\del x=y$.
\end{theorem}
\begin{proof}
Both equations hold for winner-set reset. In
\eqref{eq:absorb}, a resolved $x$ returns $x$ on both sides; an
ambiguous $x$ gives $x\cup y$ on both sides.
In \eqref{eq:left-comm}, a resolved $z$ returns $z$ on both sides;
an ambiguous $z$ gives $x\cup y\cup z$.

Assume the two equations and conservativity, and use the proposed
order. Reflexivity follows from conservativity, and antisymmetry is
immediate. If $x\del y=x$, absorption implies $y\del x=x$.
For transitivity suppose $x\leq y\leq z$. Left commutation gives
\[
 z=z\del(y\del x)=y\del(z\del x).
\]
If $z\del x=z$, the preceding absorption argument gives $x\del z=z$.
Otherwise conservativity gives $z\del x=x$, and the display gives
$z=y$, reducing to the assumed relation $x\leq y$.
If $x\leq y$ and $x\leq z$, left commutation gives
$y\del z=z\del y$; conservativity then makes $y,z$ comparable.
The operation is recovered from the order because selecting the left
argument forces the two-sided comparison just proved.

Conversely, \eqref{eq:forest-op} satisfies both equations.
For left commutation, if $z$ is below both $x,y$, these ancestors
are comparable and both sides return the higher one. If it is below
only one, both sides return that ancestor; otherwise they return $z$.
Absorption follows directly from the comparable and incomparable cases.
This order reconstruction in both directions is kernel-checked in
\texttt{ResetModels.lean}. Lemma~\ref{lem:forest-membership} below proves
that every forest operation satisfies all reset identities, completing
the equivalence.
\end{proof}

The qualification ``conservative'' is essential: the two equations are
a complete presentation within that class, not a proposed equational
basis over arbitrary algebras.

\begin{lemma}[Closure constructions]\label{lem:closures}
The reset variety is closed under the following constructions:
\begin{enumerate}
\item adjoin one new element absorbing on both sides;
\item take a disjoint union of models, using the component operation
within a component and returning the right argument across components.
\end{enumerate}
\end{lemma}
\begin{proof}
Every reset identity has the same variables on both sides: on all
ambiguous inputs its value is the union of all occurring inputs, and
marking one input by $\{a,b,c\}$ among $\{a,b\}$ detects occurrence.
Thus adjoining an absorbing element preserves every identity: either
both sides encounter it or neither does.

For the disjoint union, an evaluated term belongs to the component of
its rightmost variable. Both sides of a reset identity have the same
rightmost variable, by their behavior on atomic valuations. Fix the
set $B$ of input positions whose values belong to that component.
Prune a term with root in $B$ as follows. Retain its right child and
prune recursively; retain its left child exactly when that child's
root is in $B$. The result uses variables in $B$ and evaluates the
original term in the chosen component.

For every $C\subseteq B$, its selector is $E_t(C)\cap B$.
To see this by induction, note that $\rho(s)\notin C$ implies
$E_s(C)=\{\rho(s)\}$. Thus a removed left child contributes no
position in $B$ to the original selector, and a retained one obeys the
selector recursion. Equal original selectors give equal pruned selectors,
so the pruned terms form a reset identity. The component model satisfies
that identity. This proves closure.
\end{proof}

\begin{lemma}[Forest membership]\label{lem:forest-membership}
Every forest-order operation belongs to the reset variety.
\end{lemma}
\begin{proof}
Every finite rooted forest is constructed from singletons by disjoint
union and by adjoining an absorbing root to each component. These are
exactly the operations of Lemma~\ref{lem:closures}.
For an infinite forest order, any finite set of elements is a subalgebra,
by conservativity, and its induced order is a finite forest order.
Every evaluation of an identity lies in such a finite subalgebra.
\end{proof}

As a counting consequence, there are $(n+1)^{n-1}$ conservative reset
algebras on a fixed labeled $n$-element carrier. Adjoining a new common
root gives a bijection from rooted forests on that carrier to labeled
trees on $n+1$ vertices, so the formula is the classical tree count.
The first five counts are $1,3,16,125,1296$.

\section{Large irreducible models from combs}

Three-element completeness concerns identities. It does not bound the
irreducible models of those identities. The forest description makes
this distinction explicit.

Let $L$ be any linear order with least element $0$. Its comb has
elements
\[
 a,\qquad b_i, p_i\quad(i\in L).
\]
The $p_i$ form the ascending spine; $b_i$ is a leaf below $p_i$,
and $a$ is below every spine element. Precisely,
\[
 p_i\leq p_j\Longleftrightarrow i\leq j,\qquad
 b_i\leq p_j\Longleftrightarrow i\leq j,\qquad a\leq p_j.
\]
All other distinct pairs are incomparable. This is a forest order.
The leaves act by right projection, and its only other products are
\[
 p_i\del a=p_i,\quad
 p_i\del b_j=
 \begin{cases}p_i,&j\leq i,\\b_j,&j>i,\end{cases}\qquad
 p_i\del p_j=p_{\max(i,j)}.
\]

\begin{theorem}[Comb monolith]\label{thm:comb}
The comb is subdirectly irreducible. Its least nonidentity congruence
has the single nonsingleton class $\{a,b_0\}$.
\end{theorem}
\begin{proof}
The map replacing $a$ by $b_0$ and fixing every other element is an
endomorphism: every spine has the same product with $a$ and $b_0$.
Its kernel is the stated two-point congruence.

Any congruence identifying a spine $p_i$ with a leaf identifies $a,b_0$:
multiply the pair on the right first by $a$, then by $b_0$.
The spine returns $p_i$ in both cases, while the leaf returns $a,b_0$.
It remains to reduce each other distinct pair to a spine--leaf pair.
\begin{itemize}
\item For $a\sim b_j$ with $j>0$, multiply on the left by $p_0$ to
obtain $p_0\sim b_j$. For $j=0$ the desired pair is already identified.
\item For $b_i\sim b_j$ with $i<j$, multiply on the left by $p_i$ to
obtain $p_i\sim b_j$.
\item For $p_i\sim p_j$ with $i<j$, multiply on the right by $b_j$
to obtain $b_j\sim p_j$.
\end{itemize}
Every nonidentity congruence therefore contains the displayed kernel.
The argument, including the endomorphism and leastness, is
kernel-checked for arbitrary $L$ in \texttt{ResetModels.lean}.
\end{proof}

\begin{corollary}[Residual largeness]\label{cor:residual}
The reset variety, equivalently $\operatorname{V}(G_3^{\rm op})$, has
subdirectly irreducible models of size $2n+1$ for every $n\geq1$, and
of every infinite cardinality. In particular, it has no cardinal bound
on the sizes of its subdirectly irreducible members.
\end{corollary}
\begin{proof}
Use an $n$-element chain for the finite examples. For an infinite
cardinal $\kappa$, use its initial ordinal as $L$; the comb has size
$1+2\kappa=\kappa$. Lemma~\ref{lem:forest-membership} places every comb
in the variety, and Theorem~\ref{thm:comb} proves irreducibility.
\end{proof}

Thus a finite residual-bound argument cannot establish a finite basis
for reset. Residual largeness itself does not decide finite
axiomatizability. This conclusion concerns the reset reduct; adding
the other owned-choice operations changes congruences and does not
inherit this residual conclusion automatically.

\section{Mechanization, scope, and attribution}

\texttt{Reset.lean} proves the root ambiguity theorem, exact
selector membership semantics on mixed inputs, independence of a
position's selection from its own status, and generic completeness of
the three-element chain. It also compiles reset terms into the existing
owned-choice syntax and proves semantic correctness on nonempty input
valuations. This connects the fragment to the earlier kernel-checked
SPE interpretation. The three-element theorem applies to arbitrary input
types, without a fixed arity or a finite-case computation.

\texttt{ResetModels.lean} proves order reconstruction from the two
conservative laws, the converse forest laws, and recovery of the order.
It constructs comb forest orders, proves the collapse endomorphism,
and proves that its kernel is contained in every nonidentity congruence,
for arbitrary linearly ordered level types with a least element.
Forest membership in the reset variety, and hence the complete residual
largeness conclusion, uses the written closure-construction proof.
\texttt{OneConstant.lean} proves the constant-or-inflationary dichotomy
and its one-test rigidity consequence. The exact binary-fragment theorem
also uses the earlier mathematical atomic-projection result and has not
been fully mechanized.

The antimatroid representation, the canonical minimal-antichain
theorem, polynomial synthesis, and clone-minimality proof also have
mathematical arguments rather than complete Lean developments.
All Lean modules use 4.34.1 and only \texttt{Std}. Axiom audits report
only \texttt{propext}, \texttt{Classical.choice}, and \texttt{Quot.sound}.

\texttt{reset.py} implements normal forms, realizability, synthesis,
substitution, antimatroid products, and explicit counterexamples.
\texttt{check\_reset.py} compares independent term closure with exhaustive
antimatroid enumeration through arity four. It checks all $140^2$
four-input products. At arity five it independently examines all
$20^4$ candidate tables with one prescribed root, accepts $596$, and
constructs actual syntax for every one, giving $2980$ after varying the
root. The local DSC criterion is independently compared with greedy
constructibility on all those candidate tables. Mixed winner-set evaluations, all pairwise four-input
counterexamples, and randomized substitutions also pass. These finite
checks support, but do not replace, the generic proofs.

\texttt{forest.py} implements forest models and an independent membership
check by lifting evaluation through the full free reset algebra.
\texttt{check\_forest.py} compares every conservative operation table
through four elements with independent parent-map forest enumeration.
All $125$ four-element models lift through the full $140$-element free
algebra. The five-element comb lifts through the full $2980$-element
free algebra. Every nontrivial principal congruence is checked on combs
through fifteen elements, independently confirming the exact monolith.
The same script exhausts the single-constant unary closures for
$K=3,4,5,6$, obtaining $6,10,18,34$ functions and verifying rigidity.

The construction uses the established algebra $G_3$ and its minimal
clone, antimatroids and their basic words, DSCs and their antimatroid
correspondence, and antichain representations of monotone Boolean
functions. The exact reduction from owned choice to $G_3^{\rm op}$,
the description of its term clone and substitution by these objects,
and the forest-model and comb results are candidate contributions;
their priority remains unestablished.

This fragment loses dependence on the number of roles beyond three.
The complete owned-choice calculus remains larger and retains
owner-sensitive behavior. A finite equational basis, for reset or for
that larger calculus, has not been proved here. The exact reset normal
forms give a substantially smaller concrete target for that question.

\bibliographystyle{plain}
\bibliography{references}
\end{document}
